%% ============================================================
%%  chapitre8.tex — الفصل الثامن: دراسة الحالة وتحليل أدوار الفاعلين
%% ============================================================

\chapter{دراسة الحالة وتحليل أدوار الفاعلين}

\section{عرض الوضعية المهنية للمجموعة 8}

\begin{tcolorbox}[colback=lightgray, colframe=maincolor, arc=4pt,
    title={\textbf{وصف الوضعية المهنية}}]
يعمل السيد \textbf{أحمد بنعلي}، البالغ من العمر \textenglish{34} سنة، في مصنع \textbf{شركة النسيج المتحد} الكائن بالمنطقة الصناعية عين السبع بمدينة الدار البيضاء، بوصفه \textbf{عاملاً على الآلات} (\textenglish{opérateur de machines}) في قسم الإنتاج. وقد أُوفِد إلى هذا المصنع من طرف \textbf{وكالة الأمل للتشغيل المؤقت}، وهي مقاولة للتشغيل المؤقت مقرّها بالدار البيضاء، بموجب عقد شغل مؤقت أُبرم شفوياً دون أي وثيقة مكتوبة.

يتقاضى السيد أحمد بنعلي أجراً شهرياً قدره \textenglish{2\,800} درهم، في حين أن العمال الدائمين الذين يشتغلون في نفس المنصب ويؤدون المهام ذاتها يتقاضون \textenglish{4\,200} درهم شهرياً، أي بفارق يبلغ \textenglish{1\,400} درهم شهرياً، مما يشكّل تمييزاً صارخاً في الأجر.

امتدت فترة اشتغال السيد أحمد بنعلي بصفة متواصلة ودون انقطاع لمدة \textbf{ثلاث سنوات كاملة} (\textenglish{36} شهراً)، رغم أن مدونة الشغل تُحدّد أقصى مدة للشغل المؤقت. وتتمثل مهامه في تشغيل آلات النسيج وصيانتها الأولية ومراقبة جودة المنتجات، وهي مهام هيكلية ودائمة تدخل في صلب النشاط العادي والمستمر للمقاولة المستعملة.
\end{tcolorbox}

\vspace{0.3cm}

\begin{tcolorbox}[colback=lightgray, colframe=accentcolor, arc=4pt,
    title={\textbf{المخالفات الجوهرية المرصودة}}]
\begin{itemize}[leftmargin=*, itemsep=6pt]
    \item \textbf{غياب التصريح بالصندوق الوطني للضمان الاجتماعي} (\textenglish{CNSS}): لم تقم وكالة الأمل للتشغيل المؤقت بأي تصريح لدى الصندوق الوطني للضمان الاجتماعي طوال مدة \textenglish{36} شهراً، مما يحرم الأجير من التغطية الصحية والتعويضات العائلية ومعاش التقاعد.
    \item \textbf{غياب التغطية الصحية}: لا يستفيد السيد أحمد بنعلي من أي تأمين صحي إجباري (\textenglish{AMO})، مما يتركه دون حماية في حال تعرضه لحادثة شغل أو مرض مهني.
    \item \textbf{الخضوع المباشر لسلطة المقاولة المستعملة}: يتلقى السيد أحمد بنعلي تعليماته وأوامره مباشرة من \textbf{مدير الإنتاج} في شركة النسيج المتحد، الذي يُحدّد له ساعات العمل، والمهام اليومية، ويمارس عليه سلطة تأديبية فعلية. أما دور وكالة الأمل فيقتصر على صرف الأجر الشهري فقط.
    \item \textbf{انعدام العقد المكتوب}: لا يتوفر الأجير على أي عقد شغل مكتوب يُحدّد حقوقه والتزاماته ومدة العقد وطبيعة العمل.
    \item \textbf{التمييز في الأجر}: يتقاضى أجراً أقل بنسبة \textenglish{33\%} من الأجراء الدائمين رغم أداء نفس المهام بنفس الكفاءة.
\end{itemize}
\end{tcolorbox}

\vspace{0.3cm}

يطرح هذا الوضع إشكاليات قانونية بالغة الخطورة تتعلق بالتطبيق الصحيح لأحكام مدونة الشغل المغربية، لا سيما الأحكام المنظِّمة للتشغيل المؤقت (المواد \textenglish{495} إلى \textenglish{506})، وأحكام علاقة التبعية (المادة \textenglish{6})، والمقتضيات المتعلقة بالمسؤولية التضامنية بين المقاولة المستعملة ومقاولة التشغيل المؤقت. كما تطرح تساؤلات حول مدى مشروعية اللجوء إلى الشغل المؤقت لتغطية حاجيات هيكلية ودائمة، وحول الحماية القانونية التي يوفّرها المشرع المغربي للأجير المؤقت في مثل هذه الوضعيات.

%% ─────────────────────────────────────────────────────────
\section{تحليل دور مدير المصنع ومسؤول شركة المناولة}

\subsection{دور مدير المصنع (المقاولة المستعملة)}

يحتل مدير مصنع شركة النسيج المتحد موقعاً محورياً في هذه الوضعية، إذ يُعدّ من الناحية الفعلية والقانونية المشغّل الحقيقي للسيد أحمد بنعلي، وذلك بالنظر إلى ممارسته الفعلية لعناصر علاقة التبعية كما حدّدتها المادة \textenglish{6} من مدونة الشغل المغربية. وتتجلى مسؤوليته القانونية في عدة مستويات:

\textbf{أولاً: ممارسة السلطة الرئاسية والتأديبية بشكل مباشر.}
تنص المادة \textenglish{6} من مدونة الشغل على أن علاقة التبعية تقوم على خضوع الأجير لسلطة المشغّل الذي يُملي عليه التعليمات ويراقب تنفيذها ويملك سلطة التأديب. وفي الوضعية المدروسة، فإن مدير الإنتاج في شركة النسيج المتحد هو الذي يحدّد للسيد أحمد مواقيت العمل والمهام اليومية، ويُصدر له التوجيهات التقنية، ويمارس عليه رقابة فعلية ومستمرة. وهذه العناصر الثلاثة (التوجيه، والرقابة، والتأديب) تشكّل الأركان الجوهرية لعلاقة التبعية التي تُرتّب صفة المشغّل، بصرف النظر عن الشكل التعاقدي المعتمد. وبالتالي، فإن مدير المصنع يتحمل مسؤولية المشغّل الفعلي بمقتضى القانون.

\textbf{ثانياً: تجاوز الحدود القانونية لمدة التشغيل المؤقت.}
تنص المادة \textenglish{500} من مدونة الشغل على أن مدة عقد الشغل المؤقت لا يمكن أن تتجاوز المدة المحدّدة قانوناً، والتي ترتبط بالسبب الذي دعا إلى اللجوء إلى التشغيل المؤقت. وقد حدّدت المادة \textenglish{500} هذه المدة في ستة أشهر قابلة للتجديد مرة واحدة بالنسبة للحالات المنصوص عليها في المادة \textenglish{496}. غير أن مدير المصنع استمر في تشغيل السيد أحمد بنعلي لمدة \textenglish{36} شهراً متواصلة، أي ما يفوق الحد الأقصى القانوني بأضعاف، مما يُشكّل مخالفة صريحة لهذا المقتضى.

\textbf{ثالثاً: استخدام التشغيل المؤقت لتلبية حاجيات دائمة.}
تُحدّد المادة \textenglish{496} من مدونة الشغل الحالات التي يجوز فيها اللجوء إلى التشغيل المؤقت حصرياً في ثلاث حالات: تعويض أجير تغيّب أو توقف عقد شغله، ومواجهة تزايد مؤقت في نشاط المقاولة، وإنجاز أشغال ذات طبيعة موسمية. ولا تندرج مهام السيد أحمد بنعلي ضمن أي من هذه الحالات، لأن تشغيل آلات النسيج ومراقبة الجودة يُعدّان من الأنشطة الهيكلية والدائمة والمستمرة للمقاولة. وبالتالي فإن مدير المصنع لجأ إلى التشغيل المؤقت بصفة تعسفية بهدف التهرب من الالتزامات المترتبة على عقد الشغل غير محدد المدة.

\textbf{رابعاً: المسؤولية التضامنية في أداء الأجر.}
تُقرّ المادة \textenglish{504} من مدونة الشغل مبدأ المسؤولية التضامنية بين المقاولة المستعملة ومقاولة التشغيل المؤقت في أداء الأجر وتوابعه والتعويض عن العطلة السنوية المؤدى عنها. وتنص المواد \textenglish{86} إلى \textenglish{90} على مبدأ المساواة في الأجر بين الأجراء الذين يؤدون شغلاً متساوي القيمة. وعليه، فإن مدير المصنع يتحمل مسؤولية تضامنية عن الفارق في الأجر البالغ \textenglish{1\,400} درهم شهرياً، وعن عدم التصريح بالأجير لدى الصندوق الوطني للضمان الاجتماعي.

\begin{tcolorbox}[colback=lightgray, colframe=accentcolor, arc=4pt,
    title={\textbf{خلاصة: مخالفات مدير المصنع}}]
\begin{itemize}[leftmargin=*, itemsep=6pt]
    \item مخالفة المادة \textenglish{6}: ممارسة سلطة رئاسية مباشرة على أجير مؤقت دون الاعتراف بصفة المشغّل.
    \item مخالفة المادة \textenglish{496}: اللجوء إلى التشغيل المؤقت لتلبية حاجيات دائمة وهيكلية.
    \item مخالفة المادة \textenglish{500}: تجاوز المدة القصوى للتشغيل المؤقت (\textenglish{36} شهراً بدل \textenglish{6} أشهر).
    \item مخالفة المادتين \textenglish{86} و\textenglish{504}: التمييز في الأجر وعدم ضمان المساواة.
    \item مخالفة المواد \textenglish{89-90}: التهرب من المسؤولية التضامنية عن حقوق الأجير.
\end{itemize}
\end{tcolorbox}

\subsection{دور مسؤول شركة المناولة (وكالة الأمل للتشغيل المؤقت)}

تتحمل وكالة الأمل للتشغيل المؤقت، بوصفها مقاولة للتشغيل المؤقت، مسؤوليات قانونية جسيمة في هذه الوضعية، إذ أخلّت بجملة من الالتزامات الأساسية المنصوص عليها في مدونة الشغل والتشريعات الاجتماعية المكملة لها:

\textbf{أولاً: عدم التصريح بالأجير لدى الصندوق الوطني للضمان الاجتماعي.}
يُلزم القانون المتعلق بالضمان الاجتماعي (ظهير \textenglish{27} يوليوز \textenglish{1972}) كل مشغّل بالتصريح بأجرائه لدى الصندوق الوطني للضمان الاجتماعي وأداء الاشتراكات المستحقة. كما تُؤكّد المواد \textenglish{135} إلى \textenglish{142} من مدونة الشغل على حق الأجير في الاستفادة من أنظمة الحماية الاجتماعية. وقد امتنعت وكالة الأمل عن التصريح بالسيد أحمد بنعلي طوال مدة \textenglish{36} شهراً، مما حرمه من:
\begin{itemize}[leftmargin=*, itemsep=4pt]
    \item التعويضات العائلية عن أبنائه.
    \item التغطية الصحية الإجبارية (\textenglish{AMO}).
    \item التعويض اليومي في حالة المرض أو حادثة الشغل.
    \item حقوقه في معاش التقاعد.
    \item التعويض عن فقدان الشغل.
\end{itemize}

\textbf{ثانياً: عدم إبرام عقد شغل مكتوب.}
تنص المادة \textenglish{499} من مدونة الشغل على أن عقد الشغل المؤقت يجب أن يكون مكتوباً وأن يتضمن بيانات إلزامية محدّدة، من بينها: سبب اللجوء إلى التشغيل المؤقت، ومدة المهمة، والأجر المتفق عليه، وبيان مقاولة التشغيل المؤقت والمقاولة المستعملة. وقد أخلّت وكالة الأمل بهذا الالتزام الجوهري بعدم تحرير أي عقد مكتوب، مما يُشكّل مخالفة صريحة للمادة \textenglish{499}، ويُضعف بشكل كبير الحماية القانونية للأجير.

\textbf{ثالثاً: عدم احترام مبدأ المساواة في الأجر.}
تنص المادة \textenglish{504} من مدونة الشغل على أن أجر الأجير المؤقت لا يمكن أن يقل عن الأجر المحدد لأجير دائم يشغل نفس المنصب ويتمتع بنفس المؤهلات في المقاولة المستعملة. وبصرف الأجر الشهري البالغ \textenglish{2\,800} درهم فقط مقابل \textenglish{4\,200} درهم للعمال الدائمين، تكون الوكالة قد خالفت مبدأ المساواة في الأجر مخالفة صريحة.

\textbf{رابعاً: عدم التوفر على الضمانات القانونية المطلوبة.}
تنص المادة \textenglish{498} من مدونة الشغل على أن مقاولة التشغيل المؤقت يجب أن تتوفر على كفالة مالية لدى صندوق الإيداع والتدبير تضمن أداء أجور الأجراء المؤقتين وتوابعها في حالة عجز المقاولة عن الأداء. كما يجب عليها الحصول على ترخيص من السلطة الحكومية المكلفة بالشغل. وفي ضوء السلوك العام لوكالة الأمل، يُطرح تساؤل جدي حول مدى توفرها على هذه الضمانات القانونية.

\begin{tcolorbox}[colback=lightgray, colframe=accentcolor, arc=4pt,
    title={\textbf{خلاصة: مخالفات وكالة الأمل للتشغيل المؤقت}}]
\begin{itemize}[leftmargin=*, itemsep=6pt]
    \item مخالفة قانون الضمان الاجتماعي والمواد \textenglish{135-142} من مدونة الشغل: عدم التصريح بالأجير لدى \textenglish{CNSS}.
    \item مخالفة المادة \textenglish{499}: عدم تحرير عقد شغل مكتوب.
    \item مخالفة المادة \textenglish{504}: عدم احترام المساواة في الأجر.
    \item مخالفة محتملة للمادة \textenglish{498}: عدم التوفر على الضمانات المالية والترخيص.
    \item مخالفة المادة \textenglish{500}: المساهمة في تجاوز المدة القانونية القصوى للتشغيل المؤقت.
\end{itemize}
\end{tcolorbox}

%% ─────────────────────────────────────────────────────────
\section{تحليل دور الأجير المؤقت ومفتش الشغل}

\subsection{موقف الأجير المؤقت ومركزه القانوني}

يتمتع السيد أحمد بنعلي، رغم هشاشة وضعيته الفعلية، بحماية قانونية واسعة بمقتضى أحكام مدونة الشغل المغربية والتشريعات الاجتماعية المكملة لها. ويمكن تحديد مركزه القانوني ومجمل حقوقه على النحو التالي:

\textbf{أولاً: الحق في إعادة تكييف العقد.}
بالنظر إلى أن مدة اشتغاله تجاوزت الحد الأقصى القانوني المنصوص عليه في المادة \textenglish{500}، وأن المهام المسندة إليه ذات طابع دائم وهيكلي لا تندرج ضمن الحالات الحصرية المحدّدة في المادة \textenglish{496}، فإن الأجير يحق له المطالبة بإعادة تكييف عقده المؤقت إلى \textbf{عقد شغل غير محدد المدة} (\textenglish{CDI}) مع المقاولة المستعملة (شركة النسيج المتحد)، وذلك استناداً إلى المادة \textenglish{17} من مدونة الشغل التي تنص على أن كل عقد محدد المدة يُبرم خارج الحالات المسموح بها قانوناً يُعتبر عقداً غير محدد المدة.

\textbf{ثانياً: الحق في المساواة في الأجر.}
يحق للسيد أحمد بنعلي بمقتضى المادة \textenglish{504} من مدونة الشغل المطالبة بأجر مساوٍ لما يتقاضاه الأجراء الدائمون الذين يشغلون نفس المنصب. وعليه، فإنه يستحق الفارق في الأجر عن كامل فترة اشتغاله، أي ما مجموعه \textenglish{1\,400} درهم × \textenglish{36} شهراً = \textenglish{50\,400} درهم.

\textbf{ثالثاً: الحق في التغطية الاجتماعية بأثر رجعي.}
يحق للأجير المطالبة بتسجيله في الصندوق الوطني للضمان الاجتماعي بأثر رجعي يشمل كامل فترة الشغل (\textenglish{36} شهراً)، مع احتساب جميع الاشتراكات المستحقة التي كان يتعين على وكالة الأمل أداءها. ويشمل ذلك التعويضات العائلية والتغطية الصحية والتعويض عن فقدان الشغل.

\textbf{رابعاً: الحق في التعويض عن الأضرار.}
في حالة إنهاء علاقة الشغل بشكل تعسفي، يحق للأجير المطالبة بالتعويضات التالية:
\begin{itemize}[leftmargin=*, itemsep=6pt]
    \item \textbf{التعويض عن الفصل التعسفي}: أجر شهر ونصف عن كل سنة أقدمية أو جزء منها (المادة \textenglish{53}).
    \item \textbf{التعويض عن مهلة الإخطار}: حسب الأقدمية ووفق المادة \textenglish{43}.
    \item \textbf{التعويض عن الضرر}: شهر ونصف من الأجر عن كل سنة عمل (المادة \textenglish{41}).
    \item \textbf{فارق الأجور المستحقة}: الفرق بين الأجر المتقاضى والأجر المستحق قانوناً عن كامل فترة العمل.
    \item \textbf{التعويض المعنوي}: عن الأضرار النفسية والاجتماعية الناجمة عن الاستغلال والحرمان من الحماية الاجتماعية.
\end{itemize}

\textbf{خامساً: العوائق العملية أمام الأجير.}
رغم هذه الحقوق القانونية الواسعة، يواجه السيد أحمد بنعلي عوائق عملية تتمثل في: خشيته من فقدان مصدر رزقه في حال رفع شكوى، وضعف معرفته بحقوقه القانونية، وغياب تمثيل نقابي يدافع عن مصالحه، وصعوبة إثبات علاقة التبعية مع المقاولة المستعملة في غياب عقد مكتوب. غير أن القانون يُيسّر عبء الإثبات على الأجير، إذ يمكنه الاستناد إلى شهادة الزملاء وإلى واقع الاشتغال المستمر داخل المصنع.

\subsection{تدخل مفتش الشغل}

يحتل مفتش الشغل مكانة محورية في حماية حقوق الأجراء ومراقبة تطبيق أحكام مدونة الشغل. وفي سياق الوضعية المدروسة، يتعين عليه اتخاذ إجراءات حازمة ومتدرجة وفق الصلاحيات المخوّلة له قانوناً:

\textbf{أولاً: الزيارة الميدانية والمعاينة.}
تنص المادة \textenglish{530} من مدونة الشغل على أن مفتشي الشغل يتمتعون بصلاحية الدخول بحرية إلى أي مؤسسة خاضعة لمراقبتهم، سواء نهاراً أو ليلاً، دون إخطار مسبق. وفي وضعية السيد أحمد بنعلي، يقوم المفتش بزيارة ميدانية لمصنع شركة النسيج المتحد، سواء بناءً على شكاية الأجير أو بمبادرة تلقائية في إطار المراقبة الدورية. وخلال هذه الزيارة، يُعاين المفتش ظروف العمل الفعلية ويتحقق من طبيعة العلاقة التي تربط الأجير بالمقاولة المستعملة.

\textbf{ثانياً: الاطلاع على الوثائق والسجلات.}
تُخوّل المادة \textenglish{532} من مدونة الشغل لمفتش الشغل صلاحية الاطلاع على جميع الوثائق التي يُلزم القانون المشغّل بمسكها، ومن بينها:
\begin{itemize}[leftmargin=*, itemsep=4pt]
    \item سجل الأجراء والبيانات الشخصية لكل أجير.
    \item عقود الشغل وعقود التشغيل المؤقت.
    \item أوراق أداء الأجور (\textenglish{bulletins de paie}).
    \item وصولات التصريح بالصندوق الوطني للضمان الاجتماعي.
    \item عقد وضع الأجير رهن إشارة المقاولة المستعملة المبرم بين وكالة الأمل وشركة النسيج.
\end{itemize}

\textbf{ثالثاً: إصدار الإنذارات والملاحظات.}
تنص المواد \textenglish{539} و\textenglish{540} من مدونة الشغل على أن مفتش الشغل يمكنه توجيه ملاحظات وإنذارات للمشغّل لحمله على الامتثال للأحكام القانونية المخالَفة. وفي هذه الوضعية، يوجّه المفتش إنذاراً مكتوباً إلى كل من شركة النسيج المتحد ووكالة الأمل للتشغيل المؤقت يُنبّههما فيه إلى المخالفات المرصودة ويمنحهما أجلاً محدّداً لتسوية وضعية الأجير.

\textbf{رابعاً: تحرير المحاضر الرسمية.}
في حال عدم امتثال المشغّلين للإنذارات، يقوم مفتش الشغل بتحرير \textbf{محاضر رسمية} (\textenglish{procès-verbaux}) يُثبت فيها المخالفات المعاينة بدقة، ويُحيلها إلى النيابة العامة لدى المحكمة المختصة وفقاً للمادتين \textenglish{546} و\textenglish{547} من مدونة الشغل. وتكتسي هذه المحاضر حجية ثبوتية قوية أمام القضاء باعتبارها صادرة عن موظف عمومي أثناء أداء مهامه.

\textbf{خامساً: الدور التصالحي.}
بالإضافة إلى دوره الرقابي والزجري، يمكن لمفتش الشغل أن يقوم بدور تصالحي بين الأجير والمشغّل، من خلال عقد جلسات استماع واقتراح حلول ودية لتسوية النزاع، كتسوية وضعية الأجير بتثبيته في منصبه وأداء الفوارق المالية المستحقة.

%% ─────────────────────────────────────────────────────────
\section{الأجوبة عن أسئلة النقاش الستة}

\subsection*{السؤال الأول: هل يجوز قانوناً اللجوء إلى الشغل المؤقت في هذه الوضعية؟}

\textbf{الجواب: لا، لا يجوز قانوناً اللجوء إلى الشغل المؤقت في هذه الوضعية.} إذ تحصر المادة \textenglish{496} من مدونة الشغل المغربية حالات اللجوء إلى التشغيل المؤقت في ثلاث حالات على سبيل الحصر: أولاً، تعويض أجير تغيّب مؤقتاً أو توقف عقد شغله ما لم يكن التوقف ناتجاً عن الإضراب؛ ثانياً، مواجهة تزايد مؤقت في نشاط المقاولة؛ ثالثاً، إنجاز أشغال ذات طبيعة موسمية. وفي وضعية السيد أحمد بنعلي، فإن مهام تشغيل آلات النسيج وصيانتها ومراقبة الجودة تُعدّ من المهام الهيكلية والدائمة التي تدخل في صلب النشاط العادي والمستمر لشركة النسيج المتحد، ولا تندرج ضمن أي من الحالات الثلاث المذكورة.

علاوة على ذلك، فإن استمرار العلاقة لمدة \textenglish{36} شهراً دون انقطاع يُثبت بما لا يدع مجالاً للشك أن الحاجة إلى هذا المنصب حاجة دائمة وليست مؤقتة. وتنص المادة \textenglish{500} على أن المدة القصوى لعقد التشغيل المؤقت هي ستة أشهر غير قابلة للتجديد إلا مرة واحدة. وبالتالي، فإن اللجوء إلى التشغيل المؤقت في هذه الوضعية يُشكّل تحايلاً على القانون واستغلالاً لنظام التشغيل المؤقت لتجنب الالتزامات المترتبة على عقد الشغل غير محدد المدة، وهو ما يُبطل العقد المؤقت ويُعيد تكييفه تلقائياً إلى عقد غير محدد المدة.

\subsection*{السؤال الثاني: ما طبيعة المسؤولية التي تثيرها الوضعية؟}

\textbf{الجواب: تثير هذه الوضعية مسؤوليات متعددة ومتشعبة تشمل المسؤولية المدنية والمسؤولية الجنائية.}

على صعيد \textbf{المسؤولية المدنية التضامنية}، تُقرّ المواد \textenglish{89} و\textenglish{90} من مدونة الشغل مبدأ المسؤولية التضامنية بين المقاول الأصلي والمقاول الفرعي في أداء الأجور وتوابعها. كما تنص المادة \textenglish{504} صراحة على أن المقاولة المستعملة ومقاولة التشغيل المؤقت مسؤولتان بالتضامن عن أداء أجر الأجير المؤقت والتعويض عن العطلة السنوية. وبالتالي، فإن كلاً من شركة النسيج المتحد ووكالة الأمل للتشغيل المؤقت تتحملان مسؤولية تضامنية عن جميع الحقوق المالية المستحقة للسيد أحمد بنعلي، بما فيها فارق الأجور والتعويضات الاجتماعية.

على صعيد \textbf{المسؤولية الجنائية}، تنص المادة \textenglish{515} من مدونة الشغل على عقوبات جنائية في حق مقاولات التشغيل المؤقت التي تمارس نشاطها دون ترخيص أو دون التقيد بالأحكام المنظّمة للتشغيل المؤقت. وتشمل هذه العقوبات الغرامات المالية والحبس في حالة العود. كما يتحمل مدير المصنع مسؤولية جنائية فردية عن تشغيل أجراء مؤقتين في مخالفة للقانون، وعن عدم التصريح بالأجير لدى الصندوق الوطني للضمان الاجتماعي الذي يُشكّل جريمة معاقباً عليها بمقتضى القانون المنظم للضمان الاجتماعي.

\subsection*{السؤال الثالث: ما موقف مفتش الشغل في هذه الوضعية؟}

\textbf{الجواب: يتعين على مفتش الشغل اتخاذ موقف حازم يتضمن عدة إجراءات متدرجة.} فبموجب صلاحياته المنصوص عليها في المواد \textenglish{530} إلى \textenglish{548} من مدونة الشغل، يجب على المفتش أولاً إجراء زيارة ميدانية لمصنع شركة النسيج المتحد للمعاينة المباشرة والاطلاع على جميع الوثائق ذات الصلة (العقود، سجلات الأجراء، أوراق الأجور، وصولات التصريح بالضمان الاجتماعي). وبعد التأكد من وجود المخالفات، يوجّه إنذاراً مكتوباً إلى كل من المقاولة المستعملة ومقاولة التشغيل المؤقت يُحدّد فيه المخالفات المرصودة ويمنحهما أجلاً معقولاً لتسويتها.

في حال عدم الامتثال للإنذار، يتعين على المفتش تحرير \textbf{محاضر مخالفات رسمية} يُسجّل فيها: تجاوز المدة القانونية للتشغيل المؤقت، واللجوء إلى التشغيل المؤقت لتلبية حاجيات دائمة، والتمييز في الأجر، وعدم التصريح بالأجير لدى \textenglish{CNSS}، وغياب العقد المكتوب. ثم يُحيل هذه المحاضر إلى النيابة العامة لدى المحكمة الابتدائية المختصة وفقاً للمادتين \textenglish{546} و\textenglish{547}. كما يجب عليه إبلاغ الصندوق الوطني للضمان الاجتماعي بعدم التصريح بالأجير لاتخاذ الإجراءات القانونية اللازمة. ويمكنه أيضاً أن يأمر بتسوية وضعية الأجير فوراً من خلال تثبيته ومنحه حقوقه كاملة.

\subsection*{السؤال الرابع: ما الحلول التي يمكن تقديمها للأجير المؤقت؟}

\textbf{الجواب: يتوفر الأجير على عدة حلول قانونية يمكنه اللجوء إليها لحماية حقوقه.}

\textbf{الحل الأول}: تقديم شكاية لدى مفتشية الشغل المختصة ترابياً، يعرض فيها وضعيته ويطلب تدخل المفتش لمعاينة المخالفات وتسوية وضعه. وهذا المسلك هو الأسهل والأقل كلفة بالنسبة للأجير.

\textbf{الحل الثاني}: رفع دعوى قضائية أمام المحكمة الابتدائية (القسم الاجتماعي) يطلب فيها إعادة تكييف عقد الشغل المؤقت إلى عقد شغل غير محدد المدة مع المقاولة المستعملة (شركة النسيج المتحد)، استناداً إلى المادة \textenglish{17} وإلى مخالفة المواد \textenglish{496} و\textenglish{500}.

\textbf{الحل الثالث}: المطالبة بفارق الأجور عن كامل فترة الشغل، أي مبلغ \textenglish{50\,400} درهم (\textenglish{1\,400} درهم × \textenglish{36} شهراً)، وذلك بناءً على مبدأ المساواة في الأجر المنصوص عليه في المادة \textenglish{504}.

\textbf{الحل الرابع}: المطالبة بتسجيله لدى الصندوق الوطني للضمان الاجتماعي بأثر رجعي يشمل كامل فترة الشغل، مع أداء جميع الاشتراكات المستحقة من طرف المشغّل (حصة المشغّل وحصة الأجير).

\textbf{الحل الخامس}: في حالة فصله تعسفياً بسبب مطالبته بحقوقه، يحق له المطالبة بالتعويض عن الفصل التعسفي والتعويض عن الإخطار والتعويض عن الأقدمية والتعويض عن الضرر المعنوي.

\subsection*{السؤال الخامس: هل يمكن إعادة تكييف عقد المؤقت إلى عقد دائم؟}

\textbf{الجواب: نعم، يمكن ذلك بل إن إعادة التكييف تتم بقوة القانون في هذه الوضعية.} تنص المادة \textenglish{17} من مدونة الشغل المغربية على أنه «يُعتبر عقد الشغل المبرم لمدة محددة بمثابة عقد شغل غير محدد المدة عندما يستمر العمل به بعد أجله». وينطبق هذا المبدأ من باب أولى على عقد التشغيل المؤقت الذي يتجاوز المدة القانونية القصوى المنصوص عليها في المادة \textenglish{500}. وفي وضعية السيد أحمد بنعلي، فإن استمرار العلاقة لمدة \textenglish{36} شهراً (بدل \textenglish{6} أشهر كحد أقصى) يُحقّق شرط إعادة التكييف بصفة تلقائية.

وقد أكّدت محكمة النقض المغربية في اجتهاداتها المستقرة هذا التوجه، حيث قضت بأن الأجير المؤقت الذي يستمر في أداء عمله بعد انقضاء مدة عقده المؤقت يُعتبر مرتبطاً بعقد شغل غير محدد المدة مع المقاولة المستعملة. كما قضت بأن اللجوء إلى التشغيل المؤقت لتلبية حاجيات دائمة وهيكلية يُشكّل تحايلاً على القانون يُبرّر إعادة تكييف العلاقة. وتترتب على إعادة التكييف نتائج جوهرية أبرزها: اعتبار المقاولة المستعملة هي المشغّل الحقيقي والوحيد، واحتساب كامل فترة الشغل (\textenglish{36} شهراً) في الأقدمية، واستفادة الأجير من جميع حقوق الأجراء الدائمين بما فيها الحماية من الفصل التعسفي والحق في التعويض عن الأقدمية.

\subsection*{السؤال السادس: ما التوصيات العملية الموجّهة لأطراف الوضعية؟}

\textbf{الجواب: تتوزع التوصيات العملية على مختلف أطراف الوضعية:}

\textbf{بالنسبة لشركة النسيج المتحد (المقاولة المستعملة):} يتعين عليها المبادرة فوراً إلى تسوية وضعية السيد أحمد بنعلي بتثبيته في منصبه كأجير دائم وإبرام عقد شغل غير محدد المدة معه، وأداء فارق الأجور المستحقة عن كامل فترة الشغل، والتصريح به لدى الصندوق الوطني للضمان الاجتماعي بأثر رجعي. كما يتعين عليها مراجعة سياسة تدبير الموارد البشرية لديها والتوقف عن استخدام التشغيل المؤقت لتغطية الحاجيات الدائمة.

\textbf{بالنسبة لوكالة الأمل للتشغيل المؤقت:} يجب عليها الامتثال الفوري والكامل لأحكام المواد \textenglish{495} إلى \textenglish{506} من مدونة الشغل، بما في ذلك تحرير عقود مكتوبة لجميع الأجراء المؤقتين والتصريح بهم لدى الضمان الاجتماعي واحترام مبدأ المساواة في الأجر. وفي حال استمرارها في المخالفة، فإنها تُعرّض نفسها لسحب الترخيص والعقوبات الجنائية المنصوص عليها في المادة \textenglish{515}.

\textbf{بالنسبة للأجير (السيد أحمد بنعلي):} يُنصح بتوثيق كل ما يُثبت علاقته بشركة النسيج المتحد (شهادات الزملاء، جداول الحضور، أي وثيقة تحمل اسمه صادرة عن المصنع)، والتوجه إلى مفتشية الشغل المختصة لتقديم شكايته، والاستعانة بمحامٍ متخصص في قانون الشغل لرفع دعوى إعادة التكييف والمطالبة بحقوقه كاملة. كما يُنصح بالتواصل مع النقابة المهنية للقطاع للحصول على الدعم والمساندة.

\textbf{بالنسبة لمفتش الشغل:} يتعين عليه تطبيق القانون بصرامة وحزم، وعدم الاكتفاء بالإنذارات الشفوية بل تحرير محاضر رسمية وإحالتها إلى النيابة العامة. كما يُوصى بتكثيف الزيارات الميدانية الرقابية للمقاولات التي تلجأ إلى التشغيل المؤقت، وبالتنسيق مع الصندوق الوطني للضمان الاجتماعي لمحاربة ظاهرة عدم التصريح بالأجراء.

\textbf{على صعيد التوصيات العامة والإصلاح التشريعي:} يُوصى بتشديد العقوبات المقررة لمخالفة أحكام التشغيل المؤقت لتكون رادعة فعلاً، وبإنشاء آلية رقابية رقمية مركزية تُتيح تتبع عقود التشغيل المؤقت ومدتها في الزمن الحقيقي، وبتعزيز أعداد مفتشي الشغل وتوفير الإمكانيات اللوجستية لهم لأداء مهامهم الرقابية بفعالية، وبتوعية الأجراء المؤقتين بحقوقهم القانونية من خلال حملات تحسيسية تنظمها وزارة الإدماج الاقتصادي والمقاولة الصغرى والتشغيل والكفاءات بالتعاون مع النقابات المهنية.
